% Generated draft: requires mathematical and source review.
\documentclass[11pt]{article}
\usepackage[margin=25mm]{geometry}
\usepackage{fontspec}
\setmainfont{Latin Modern Roman}
\usepackage{amsmath,amssymb,mathtools,unicode-math}
\setmathfont{Latin Modern Math}
\usepackage{enumitem,needspace,xurl,hyperref}
\setlength{\parindent}{0pt}
\setlength{\parskip}{5pt}
\setlength{\emergencystretch}{3em}
\newcommand{\N}{\mathbb{N}}
\newcommand{\Z}{\mathbb{Z}}
\newcommand{\Q}{\mathbb{Q}}
\newcommand{\R}{\mathbb{R}}
\newcommand{\C}{\mathbb{C}}
\newcommand{\F}{\mathbb{F}}
\newcommand{\A}{\mathbb{A}}
\newcommand{\PP}{\mathbb{P}}
\newcommand{\E}{\mathbb{E}}
\newcommand{\eps}{\varepsilon}
\newcommand{\dd}{\,\mathrm{d}}
\begin{document}
\section*{UnsolvedMath ID: 30001405}
\textbf{Source problem code:} catalogue entry
\subsection{Definitions and mathematical statement}
Let \(X\) be definable in an o-minimal structure, and let \(E\) be a type-definable equivalence relation on \(X\) of bounded index. Here type-definable means expressible as an intersection of definable sets; bounded index is understood relative to the saturation cardinal of the ambient model. The supplied record does not specify that ambient model or its cardinal conventions. Write \(q\colon X\to X/E\) for the quotient map and equip \(X/E\) with the logic topology: a subset \(C\subseteq X/E\) is closed when \(q^{-1}(C)\) is type-definable, using the ambient model and parameter conventions of the quotient construction.

Assume that \(X/E\) is locally contractible and that, for every \(x\in X\), there are definable sets \(D_{x,k}\), indexed by \(k\in\N\), such that
\[
D_{x,k+1}\subseteq D_{x,k},\qquad
[x]_E=\bigcap_{k\in\N}D_{x,k},
\]
and every \(D_{x,k}\) is definably contractible. Definable contractibility means that the identity map is homotopic to a constant map through a definable homotopy. No uniform definability of the families as \(x\) varies is asserted in the supplied statement. The record does not specify whether the sets \(D_{x,k}\) must lie in \(X\), nor does it specify a precise convention for local contractibility.

For a basepoint \(x_0\in X\), write \(\pi_n^{\mathrm{def}}(X,x_0)\) for the o-minimal homotopy group formed using based definable maps and definable homotopies, and write \(\pi_n(X/E,q(x_0))\) for the ordinary based homotopy group of the logic-topological quotient. The conjectural comparison requested by the expanded source statement is
\[
\pi_n^{\mathrm{def}}(X,x_0)\cong
\pi_n(X/E,q(x_0)).
\]
The source asks for this comparison ``for every \(n\).'' A based group formulation asks whether it holds for every positive integer \(n\) and every choice of \(x_0\in X\); however, the original wording refers to \emph{higher} homotopy groups, which may mean only \(n\geq2\). If degree zero is intended, it requires a separate statement about component sets rather than group isomorphisms. These degree and basepoint conventions remain unresolved.

The original formulation invokes strong assumptions labelled \((\mathrm{A1}^{*})\) and \((\mathrm{A2}^{*})\). Supplied verification metadata says that the displayed contractibility hypotheses expand those assumptions by replacing simply connected hypotheses with contractibility, but their complete original definitions are not supplied. Thus equivalence with the original labelled assumptions is not independently established here. The precise ambient o-minimal homotopy framework and whether a particular comparison map or natural isomorphism is intended also require clarification.
\subsection{Short English statement}
The conjecture asks whether the o-minimal homotopy groups of a definable set agree with the ordinary homotopy groups of its logic-topological quotient when the quotient is locally contractible and each equivalence class is a countable decreasing intersection of definably contractible sets. The exact strong assumptions, ambient framework, degree range, and basepoint conventions require review. The supplied record labels the problem open; this draft does not independently verify that status.
\subsection{Sources}
\begin{itemize}
\item Local registry: \texttt{data/unsolvedmath/problems.json}, permanent UnsolvedMath ID \(30001405\), title \emph{Homotopy Groups of Definable Quotients}; source problem code \texttt{catalogue entry}. The supplied record belongs to the Oberwolfach Reports Open Problems collection and gives the proposed year as 2010; no individual proposer is identified.
\item Supplied report citation: \emph{Model Theory: Around Valued Fields and Dependent Theories} (2010), Oberwolfach Reports, pp.~5--53, DOI \texttt{10.4171/owr/2010/01}. Source URL and supplied literature-source link labelled ``original Oberwolfach problem report'': \url{https://doi.org/10.4171/owr/2010/01}.
\item Supplied literature assessment, checked on \texttt{2026-08-21}: the passage explicitly presents the higher-homotopy identification as a conjecture, and the supplied curation reports locating no primary paper establishing it in the stated generality. Its listed checked source is the original report linked above. This is supplied dated curation, not an independently verified literature assessment; time-sensitive status and cited claims require independent verification.
\item Supplied statement metadata labels the formulation \texttt{corrected\_verified} and states that the official report defines the strong assumptions by replacing simply connected hypotheses with contractibility. These are supplied verification claims; this draft does not claim to have consulted the report.
\end{itemize}
\end{document}
